\documentclass[11pt,fleqn]{report}
\usepackage{graphicx}
\usepackage{amssymb}
\usepackage{longtable}
\usepackage[T1]{fontenc}

\usepackage{cml}
\lstnewenvironment{codeXML}
    {\lstset{basicstyle=\small\ttfamily,
             language=XML,
             xleftmargin=0.0in,
             xrightmargin=0.0in,
             backgroundcolor=\color{Orchid!60}, % 60% orchid, 40% white
             escapeinside={\%*}{*)}, %  .tex file characters inside %*...*) are
                                     %  treated as Tex commands.
            }
    }
    { }

\newcommand{\ModelName}{Convert String}
\newcommand{\ModelAuthor}{Daniel Ghan \\ Gary Turner}
\newcommand{\ModelAffiliation}{Odyssey Space Research}
\newcommand{\ModelDate}{July 2022}
\pagestyle{plain}
\usepackage[parfill]{parskip}
\usepackage[margin=1.5cm]{caption}
\usepackage{subfig}

\begin{document}
\titlePage
\clearpage % Reset page number
\pagenumbering{roman}
\tableofcontents % Print table of contents
\vfill
{\Large \textbf {Revision History}}
{\large
\begin{center}
\begin{tabular}{|l||l|l|l|}
\hline
 \textbf{Version} & \textbf{Date} & \textbf{Author} & \textbf{Purpose} \\ \hline
 \hline
1 & July 2022 & Gary Turner & Initial Version \\ \hline
\end{tabular}
\end{center}
}

\chapter{Introduction} % Make a "chapter" heading
\pagenumbering{arabic} % Start text with arabic 1
The Convert String model supports the conversion of numeric data read in
as a string, or char-array into a numeric data type. This model was written
to support the fault-injection / fault-management CML model which parses
XML data and generates numeric data from those data.


\chapter {Requirements}
The model shall support conversion to the following data types:
\begin{itemize}
  \item float
  \item double
  \item int
  \item unsigned int
  \item long
  \item unsigned long
  \item long long
  \item unsigned long long
  \item size\_t
  \item bool
\end{itemize}


\chapter {Model Specification}
The model is written to provide a single user-interface templated method,
convert which can take either a STL-string or \textit{char} pointer as its
single argument and returns the converted value.

For example:
\begin{codeC}
double d = ConvertString::convert<double>("12345.6789")
\end{codeC}

or, more commonly:
\begin{codeC}
std::string s = "12345.6789"
double d = ConvertString::convert<double>(s)
\end{codeC}
Note that no effort has been made to support or evaluate SWIG compatibility
(e.g. for calling this method from a Trick input file).
There is no known use-case at this time to require such support.


\chapter {User's Guide}
The model provides only the \textit{convert}
method presented in the previous section.
\chapter{Verification}

\section{Code Coverage}
\begin{figure}[h]
  \centering
  \includegraphics[scale=0.4]{graphics/code_coverage.png}
  \caption{Code Coverage}
  \label{fig:coverage}
\end{figure}
Note -- there is a sanity check in
\textit{ConvertString::convert(const char* str)} that prevents calls going
to \textit{convert\_numeric<T>} in the case that $T$ is a non-arithmetic type.
However, to the extent of verification testing, there appears to be no type
$T$ that is both non-arithmetic and allows model compilation;
see notes in the header file to explain this.

This is a template function so for every available $T$, the test
\textit{is\_arithmetic<T>} passes unambiguously and the compiler considers
only the path following success of the sanity check.

Thus, the escape
path and error message that would be executed upon failing the sanity check
\textit{does not even get compiled}.

Consequently, this code is unreachable; indeed, it is
\textit{non-compilable}. It
does not show up as non-executed code in the code-coverage because it is
not compiled into the executive, and is bypassed as a comment would be.

\newpage
\section {Verification Tests}
The model provides a single test case, at \textit{SIM\_verif/RUN\_verif}.
This test case exercises conversion to all of the required data types.
The Boolean type is tested with "true", "FALSE", and a numerical
conversion which generates an error message as an invalid argument,
and converted to false. An additional error test case investigates
the detection of a call using a NULL pointer. Both errors produce
messages as expected.
The single pass is recorded for regression testing in
\textit{verif\_data/RUN\_verif}.
\begin{figure}[h]
  \centering
  \includegraphics[scale=0.5]{graphics/convert_string_verif.PNG}
\end{figure}

Note that the value for \textit{test.f} (a float) is truncated from
the value for \textit{test.d} (a double).

\end{document}
